% from ega2-1-fr.tex
\oldpage[II]{106}

La question étant locale sur $Y$, on peut supposer
$Y=\operatorname{Spec}(A)$~; en outre
$X_{f^{\otimes n}}=X_f$, donc en remplaçant $\mathcal{L}$ par un
$\mathcal{L}^{\otimes n}$ convenable, on peut supposer $\mathcal{L}$ très
ample pour le morphisme structural $q:X\to Y$
(\hyperref[II.4.6.11-fr]{4.6.11}). L'homomorphisme canonique
$\sigma:q^*(q_*(\mathcal{L}))\to\mathcal{L}$ est donc surjectif et le
morphisme correspondant
\[
r=r_{\mathcal{L},\sigma}:X\to P=\mathbf{P}(q_*(\mathcal{L}))
\]
une immersion telle que $\mathcal{L}=r^*(\mathcal{O}_P(1))$
(\hyperref[II.4.4.4-fr]{4.4.4})~; en outre, comme $X$ est propre sur $Y$,
l'immersion $r$ est fermée (\hyperref[II.5.4.4-fr]{5.4.4}). Par ailleurs, on
a par définition $f\in\Gamma(Y,q_*(\mathcal{L}))$ et $\sigma^\flat$ est
l'identité de $q_*(\mathcal{L})$~; il résulte alors de la formule
(\hyperref[II.3.7.3.1-fr]{3.7.3.1}) que l'on a $X_f=r^{-1}(D_+(f))$~; $X_f$
est par suite un sous-préschéma fermé du schéma affine $D_+(f)$ et est donc
bien un schéma affine.

Lorsqu'on prend en particulier $Y=X$, on obtient (compte tenu de
(\hyperref[II.4.6.13-fr]{4.6.13, (i)})) le corollaire suivant, dont la
démonstration directe est d'ailleurs immédiate~:

\begin{corollary}[5.5.8]
\label{II.5.5.8-fr}
Soit $X$ un préschéma, $\mathcal{L}$ un $\mathcal{O}_X$-Module inversible.
Pour toute section $f$ de $\mathcal{L}$ au-dessus de $X$, $X_f$ est affine
sur $X$ (et par suite un schéma affine lorsque $X$ est un schéma affine).
\end{corollary}

\subsection{Le lemme de Chow.}
\label{subsection:II.5.6-fr}

\begin{theorem}[5.6.1]
\label{II.5.6.1-fr}
(Lemme de Chow). Soit $S$ un préschéma, $X$ un $S$-schéma de type fini. On
suppose vérifiée l'une des hypothèses suivantes~:
\begin{enumerate}
  \item[\rm{a)}] $S$ est noethérien.
  \item[\rm{b)}] $S$ est un schéma quasi-compact, et $X$ a un nombre fini de
  composantes irréductibles.
\end{enumerate}
Dans ces conditions~:
\begin{enumerate}
  \item[\rm{(i)}] Il existe un $S$-schéma $X'$ quasi-projectif sur $S$ et un
  $S$-morphisme $f:X'\to X$, projectif et surjectif.
  \item[\rm{(ii)}] On peut prendre $X'$ et $f$ tels qu'il existe un ouvert
  $U\subset X$ pour lequel $U'=f^{-1}(U)$ soit dense dans $X'$ et la
  restriction de $f$ à $U'$ un isomorphisme $U'\simeq U$.
  \item[\rm{(iii)}] Si $X$ est réduit (resp. irréductible, intègre), on peut
  supposer $X'$ réduit (resp. irréductible, intègre).
\end{enumerate}
\end{theorem}

La démonstration se fait en plusieurs étapes.

A) On peut d'abord se ramener au cas où $X$ est irréductible. En effet, dans
l'hypothèse a), $X$ est noethérien, donc, dans les deux hypothèses, les
composantes irréductibles $X_i$ de $X$ sont en nombre fini. Si le théorème est
démontré pour chacun des sous-préschémas fermés réduits de $X$ ayant les $X_i$
pour espaces sous-jacents, et si $X'_i$ et $f_i:X'_i\to X_i$ sont le préschéma
et le morphisme correspondant à $X_i$, le préschéma $X'$ somme des $X'_i$, et
le morphisme $f:X'\to X$ dont la restriction à chaque $X'_i$ est $j_i\circ f_i$
($j_i$ injection canonique $X_i\to X$), répondent à la question. Il est
immédiat en effet que $X'$ est réduit si les $X'_i$ le sont~; d'autre part on
satisfait à (ii) en prenant pour $U$ la réunion des ensembles
$U_i\cap(\bigcup_{j\neq i}X_j)^{\complement}$. Enfin, comme les $X'_i$ sont
quasi-projectifs sur $S$, il en est de même de $X'$

